\chapter{Line-by-line: \texttt{piret\_sm4.py}}
\section{\texttt{sm4\_fault\_encrypt}}
Simulate a controlled fault during the real SM4 implementation, while collecting teacher-only intermediates.
\begin{center}\begin{tikzpicture}[node distance=4mm]
\node[box,text width=100mm] (input) {Plaintext and target key};
\node[box,below=of input,text width=100mm] (action) {Run SM4; replace selected S-box input};
\node[box,below=of action,text width=100mm] (output) {Faulty ciphertext and teaching trace};
\draw[flow](input)--(action);\draw[flow](action)--(output);
\end{tikzpicture}\end{center}
\remember{\textbf{Physical interpretation:} Target-side behavior is modeled here. In firmware, intermediate words could occupy CPU registers or RAM depending on compilation; this Python callback does not access STM32 registers.}
\subsection*{Exact source}
\VerbatimInput[firstnumber=14]{generated/annotated_piret_sm4_sm4_fault_encrypt.txt}
\subsection*{Read each executable line}
\begingroup\small\begin{longtable}{r >{\raggedright\arraybackslash}p{138mm}}\toprule Line&Meaning and effect\\\midrule\endfirsthead\toprule Line&Meaning and effect\\\midrule\endhead
14&Declare this callable and its parameters; the body runs only when called. Simulate a controlled fault during the real SM4 implementation, while collecting teacher-only intermediates.\\
21&If a fault was requested, require a byte position 0..3 and a nonzero one-byte mask. None permits a clean reference run.\\
22&Stop this path with the stated exception. The message identifies the violated input, model or verification condition; no recovered key is accepted by this line.\\
24&Reject a selected round outside the zero-based 0..31 interval.\\
25&Stop this path with the stated exception. The message identifies the violated input, model or verification condition; no recovered key is accepted by this line.\\
28&Compute the target schedule for trace labeling. The result is a list indexed 0..31, not a dictionary; it is not passed to the output-only solver.\\
31&Create an empty teaching trace; the following explicit loop adds one record for each round.\\
33&Visit all thirty-two round indices to construct the trace records.\\
34&Add the record for this round and its target-side round key to the teacher trace.\\
36&Translate cipher callback event names into short trace dictionary fields a, b, linear and new.\\
38&Define the target callback; its body executes later when the normal cipher calls the observation hook, not when this definition is reached.\\
41&Inject only when both round and operation match, and a target byte was selected; other events pass through unchanged.\\
42&Move the mask into the selected big-endian byte and XOR it into the word. For byte 0 and mask 01, the shifted mask is 01000000.\\
44&Record only named scalar intermediate events; ignore other observation events for this teacher trace.\\
45&Save the observed scalar at its round and semantic field, so the terminal can compare clean and faulty calculations.\\
47&Return the unchanged or injected scalar to the normal SM4 engine; this returned replacement is what the next cipher operation consumes.\\
49&Execute normal SM4 with the external observer attached. A returned observer replacement is consumed by the real round function.\\
51&Return the target output together with teacher-only intermediate records. The oracle later forwards only the ciphertext to the solver.\\
\bottomrule\end{longtable}\endgroup
\section{\texttt{inverse\_rows}}
Construct and cache the inverse of the public 32-bit diffusion map.
\begin{center}\begin{tikzpicture}[node distance=4mm]
\node[box,text width=100mm] (input) {Public SM4 linear map L};
\node[box,below=of input,text width=100mm] (action) {Binary Gaussian elimination on [M | I]};
\node[box,below=of action,text width=100mm] (output) {Rows of inverse L};
\draw[flow](input)--(action);\draw[flow](action)--(output);
\end{tikzpicture}\end{center}
\remember{\textbf{Physical interpretation:} This is attacker-side algebra on the PC. No pulse, board read or secret key is required.}
\subsection*{Exact source}
\VerbatimInput[firstnumber=54]{generated/annotated_piret_sm4_inverse_rows.txt}
\subsection*{Read each executable line}
\begingroup\small\begin{longtable}{r >{\raggedright\arraybackslash}p{138mm}}\toprule Line&Meaning and effect\\\midrule\endfirsthead\toprule Line&Meaning and effect\\\midrule\endhead
54&Cache the public inverse-matrix result after the first calculation. The matrix is key-independent and is not recomputed for every observation.\\
55&Declare this callable and its parameters; the body runs only when called. Construct and cache the inverse of the public 32-bit diffusion map.\\
59&Prepare the list of images of the thirty-two input basis words under the public map L.\\
61&Choose each input basis bit explicitly, from bit zero through bit thirty-one.\\
62&Construct a 32-bit word with exactly the selected input bit set.\\
64&Apply public encryption diffusion L to that basis word, obtaining one matrix column.\\
66&Store this basis image as the next column of the binary matrix.\\
69&Prepare the augmented-matrix row list before encoding individual row bits.\\
71&Construct each output row of the binary matrix explicitly.\\
72&Start this encoded matrix row with every bit cleared.\\
74&Read the selected output bit from each stored matrix column.\\
75&Extract the bit indicating whether this input column contributes to the selected output row.\\
76&Place that contribution bit in the encoded matrix row; OR combines disjoint bit positions.\\
78&Place the matching identity-matrix bit in the upper half of the augmented row.\\
80&Combine the L-matrix row with its identity bit to form the packed row of [M | I].\\
82&Store the completed augmented row before Gaussian elimination.\\
84&Process one pivot column of the 32-by-32 binary matrix at a time.\\
86&Mark the pivot as not yet found; the following ordinary loop searches eligible rows.\\
88&Search the current and remaining rows for a one in the pivot column.\\
89&Read the current candidate row bit at the pivot column.\\
91&Accept this row as a valid pivot only if its selected column bit is one.\\
92&Remember the actual row index, then stop searching because a valid pivot was found.\\
93&Stop searching pivot rows as soon as a valid row has been found.\\
95&Detect a missing pivot instead of blindly swapping an invalid row index.\\
96&Stop this path with the stated exception. The message identifies the violated input, model or verification condition; no recovered key is accepted by this line.\\
98&Swap the chosen pivot row into the current diagonal position.\\
100&Visit every augmented row so the current column can be cleared outside its pivot.\\
101&Leave the pivot unchanged; eliminate only another row whose pivot-column bit is one.\\
103&XOR the entire augmented pivot row into that row, performing GF(2) row elimination on both halves.\\
106&Prepare the final public inverse-matrix rows after elimination.\\
108&Visit each reduced augmented row to extract its inverse half.\\
109&Discard the lower identity half and retain the upper inverse-matrix row.\\
111&Add the extracted inverse row in the original output-bit order.\\
113&Return an immutable sequence of inverse rows, cached for reuse by later observations.\\
\bottomrule\end{longtable}\endgroup
\section{\texttt{inverse\_linear}}
Apply the inverse binary matrix to one observed output difference.
\begin{center}\begin{tikzpicture}[node distance=4mm]
\node[box,text width=100mm] (input) {32-bit output XOR};
\node[box,below=of input,text width=100mm] (action) {Compute inverse-matrix row parities};
\node[box,below=of action,text width=100mm] (output) {32-bit S-box output XOR};
\draw[flow](input)--(action);\draw[flow](action)--(output);
\end{tikzpicture}\end{center}
\remember{\textbf{Physical interpretation:} The PC transforms received ciphertext bytes into a mathematical difference. The target does not run backward.}
\subsection*{Exact source}
\VerbatimInput[firstnumber=116]{generated/annotated_piret_sm4_inverse_linear.txt}
\subsection*{Read each executable line}
\begingroup\small\begin{longtable}{r >{\raggedright\arraybackslash}p{138mm}}\toprule Line&Meaning and effect\\\midrule\endfirsthead\toprule Line&Meaning and effect\\\midrule\endhead
116&Declare this callable and its parameters; the body runs only when called. Apply the inverse binary matrix to one observed output difference.\\
119&Clear the output word before placing one parity bit at a time.\\
121&Associate each inverse-matrix row with the output bit it computes.\\
122&Select exactly the input bits used by this inverse row.\\
124&Count selected one bits; only whether this count is odd or even matters over GF(2).\\
126&Keep the lowest count bit: odd becomes one and even becomes zero.\\
128&Place this row parity into its output position without changing other output bits.\\
130&Return the word assembled from all thirty-two inverse-row parity calculations.\\
\bottomrule\end{longtable}\endgroup
\section{\texttt{sm4\_byte\_candidates}}
Find key-byte guesses consistent with one restored ciphertext pair and a known byte fault.
\begin{center}\begin{tikzpicture}[node distance=4mm]
\node[box,text width=100mm] (input) {C, F, byte index and known mask};
\node[box,below=of input,text width=100mm] (action) {Undo L; test every S-box input};
\node[box,below=of action,text width=100mm] (output) {Compatible key-byte set};
\draw[flow](input)--(action);\draw[flow](action)--(output);
\end{tikzpicture}\end{center}
\remember{\textbf{Physical interpretation:} This filtering happens on the PC after a response has been received. The known-mask assumption must be established separately for a physical experiment.}
\subsection*{Exact source}
\VerbatimInput[firstnumber=133]{generated/annotated_piret_sm4_sm4_byte_candidates.txt}
\subsection*{Read each executable line}
\begingroup\small\begin{longtable}{r >{\raggedright\arraybackslash}p{138mm}}\toprule Line&Meaning and effect\\\midrule\endfirsthead\toprule Line&Meaning and effect\\\midrule\endhead
133&Declare this callable and its parameters; the body runs only when called. Find key-byte guesses consistent with one restored ciphertext pair and a known byte fault.\\
136&Require two complete block outputs and valid controlled-fault descriptors before interpreting differences.\\
137&Stop this path with the stated exception. The message identifies the violated input, model or verification condition; no recovered key is accepted by this line.\\
139&Decode each sixteen-byte output into four big-endian unsigned 32-bit words.\\
142&Require the last three restored output words to be unchanged. Otherwise this observation violates the selected last-round model.\\
143&Stop this path with the stated exception. The message identifies the violated input, model or verification condition; no recovered key is accepted by this line.\\
146&XOR the newest words, then undo public linear diffusion to recover the S-box output difference.\\
148&Map byte 0,1,2,3 to bit shifts 24,16,8,0 in a big-endian 32-bit word.\\
150&Reject a zero difference or any nonzero difference bit outside the selected byte.\\
152&Stop this path with the stated exception. The message identifies the violated input, model or verification condition; no recovered key is accepted by this line.\\
154&Select the affected S-box output difference byte, keeping only its lowest eight bits.\\
157&Compute q: select a byte of the XOR of the three observable unchanged words. It is not a hidden state read.\\
160&Start an empty set of compatible unknown S-box input bytes.\\
162&Test each possible unknown input byte explicitly.\\
163&Apply the known error hypothesis to the candidate S-box input.\\
165&Calculate the clean/faulty S-box XOR for this candidate input and known error.\\
167&Keep the hypothesis only if its predicted S-box difference equals the observed one.\\
168&Record this compatible input value; the set removes duplicates.\\
170&Prepare key-byte candidates separately from S-box input hypotheses.\\
172&Convert each compatible S-box input into a key-byte guess.\\
173&Solve a = q XOR key\_byte using the observable q byte.\\
175&Store the resulting key-byte hypothesis for intersection with other observations.\\
177&Return the key hypotheses, compatible S-box inputs and full inverse-L difference separately.\\
\bottomrule\end{longtable}\endgroup
\section{\texttt{reverse\_sm4}}
Derive the initial key words from four consecutive recovered tail words, then remove FK.
\begin{center}\begin{tikzpicture}[node distance=4mm]
\node[box,text width=100mm] (input) {Recovered rk28 through rk31};
\node[box,below=of input,text width=100mm] (action) {Reverse schedule; remove public FK};
\node[box,below=of action,text width=100mm] (output) {Original 128-bit master key};
\draw[flow](input)--(action);\draw[flow](action)--(output);
\end{tikzpicture}\end{center}
\remember{\textbf{Physical interpretation:} All schedule inversion runs on the PC. Public CK/FK constants are not extracted from the board.}
\subsection*{Exact source}
\VerbatimInput[firstnumber=180]{generated/annotated_piret_sm4_reverse_sm4.txt}
\subsection*{Read each executable line}
\begingroup\small\begin{longtable}{r >{\raggedright\arraybackslash}p{138mm}}\toprule Line&Meaning and effect\\\midrule\endfirsthead\toprule Line&Meaning and effect\\\midrule\endhead
180&Declare this callable and its parameters; the body runs only when called. Derive the initial key words from four consecutive recovered tail words, then remove FK.\\
183&Require four consecutive tail subkeys; a single recovered rk31 is insufficient to seed the inverse schedule.\\
184&Stop this path with the stated exception. The message identifies the violated input, model or verification condition; no recovered key is accepted by this line.\\
187&Reserve K0..K31 as unknown and seed K32..K35 with recovered rk28..rk31 in that order.\\
189&Recover schedule words in descending order so every right-hand dependency is already known.\\
192&Form the public schedule-core input from three known neighboring words and CK[index].\\
194&Apply the forward S-box and key-specific L-prime, then XOR with K[index+4] to solve for the older schedule word.\\
196&Call optional explanatory instrumentation only if supplied; the numerical result has already been computed.\\
197&Report this derived schedule word without giving the observer control over key recovery.\\
200&Prepare four unwhitened master-key words.\\
202&Remove FK separately for master-word positions zero through three.\\
203&Undo initial key whitening with the corresponding public FK constant.\\
204&Preserve the master-word order for final serialization.\\
206&Serialize the four recovered master words into sixteen big-endian bytes.\\
\bottomrule\end{longtable}\endgroup
\section{\texttt{peel\_sm4}}
Reconstruct an earlier output window using only already recovered later keys.
\begin{center}\begin{tikzpicture}[node distance=4mm]
\node[box,text width=100mm] (input) {Ciphertext plus recovered later keys};
\node[box,below=of input,text width=100mm] (action) {Undo later SM4 round recurrences};
\node[box,below=of action,text width=100mm] (output) {Restored target-round output view};
\draw[flow](input)--(action);\draw[flow](action)--(output);
\end{tikzpicture}\end{center}
\remember{\textbf{Physical interpretation:} This is offline PC computation. It neither rewinds the STM32 nor reads an old state from target RAM.}
\subsection*{Exact source}
\VerbatimInput[firstnumber=209]{generated/annotated_piret_sm4_peel_sm4.txt}
\subsection*{Read each executable line}
\begingroup\small\begin{longtable}{r >{\raggedright\arraybackslash}p{138mm}}\toprule Line&Meaning and effect\\\midrule\endfirsthead\toprule Line&Meaning and effect\\\midrule\endhead
209&Declare this callable and its parameters; the body runs only when called. Reconstruct an earlier output window using only already recovered later keys.\\
217&Reuse a sixteen-byte length check. This does not execute AES and does not import an AES key into SM4 recovery.\\
219&Restrict this full-recovery peeling helper to the four modeled target rounds.\\
220&Stop this path with the stated exception. The message identifies the violated input, model or verification condition; no recovered key is accepted by this line.\\
223&Read the four final output words, then reverse their order to restore the chronological rolling state.\\
225&Undo only rounds later than the requested target round, newest first.\\
226&Require a genuinely recovered later key before reversing its round; never substitute a reference key.\\
227&Stop this path with the stated exception. The message identifies the violated input, model or verification condition; no recovered key is accepted by this line.\\
231&Recompute the later-round nonlinear/linear transform and XOR with its newest word to reconstruct the missing oldest word.\\
233&Shift the reconstructed window backward by one round: old, u0, u1, u2.\\
236&Return a synthetic ciphertext-style view with word reversal, allowing reuse of the same last-round candidate equation.\\
\bottomrule\end{longtable}\endgroup
\section{\texttt{recover\_sm4\_master\_key}}
Coordinate independent fault observations, intersections, four tail-key recoveries and inverse key expansion.
\begin{center}\begin{tikzpicture}[node distance=4mm]
\node[box,text width=100mm] (input) {Correct output and fault oracle};
\node[box,below=of input,text width=100mm] (action) {Recover rk31, rk30, rk29, rk28};
\node[box,below=of action,text width=100mm] (output) {Master key or explicit ambiguity};
\draw[flow](input)--(action);\draw[flow](action)--(output);
\end{tikzpicture}\end{center}
\remember{\textbf{Physical interpretation:} On hardware, each oracle request would involve another target encryption and response. The implementation here makes software calls, not USB or serial requests.}
\subsection*{Exact source}
\VerbatimInput[firstnumber=239]{generated/annotated_piret_sm4_recover_sm4_master_key.txt}
\subsection*{Read each executable line}
\begingroup\small\begin{longtable}{r >{\raggedright\arraybackslash}p{138mm}}\toprule Line&Meaning and effect\\\midrule\endfirsthead\toprule Line&Meaning and effect\\\midrule\endhead
239&Declare this callable and its parameters; the body runs only when called. Coordinate independent fault observations, intersections, four tail-key recoveries and inverse key expansion.\\
248&Reuse a sixteen-byte length check. This does not execute AES and does not import an AES key into SM4 recovery.\\
250&Freeze a possibly one-shot iterable so every target byte can reuse the same ordered observation masks.\\
252&Reject an empty fault-mask sequence before any target query.\\
253&Stop this path with the stated exception. The message identifies the violated input, model or verification condition; no recovered key is accepted by this line.\\
255&Reject repeated masks; deterministic repeated observations do not add constraints.\\
256&Stop this path with the stated exception. The message identifies the violated input, model or verification condition; no recovered key is accepted by this line.\\
258&Validate each supplied fault descriptor explicitly.\\
259&Reject zero and masks that do not fit in one nonzero byte.\\
260&Stop this path with the stated exception. The message identifies the violated input, model or verification condition; no recovered key is accepted by this line.\\
262&Require the first selected byte position to be zero through three.\\
263&Stop this path with the stated exception. The message identifies the violated input, model or verification condition; no recovered key is accepted by this line.\\
265&Initialize recovered keys, event history, per-byte counts and actual faulty-output query count separately.\\
267&Target rounds 31,30,29,28 in that order; later recovered keys enable peeling for the next target.\\
269&Restore the clean reference to this target round using only the currently recovered later keys.\\
271&Allocate four output bytes for the current round key; unfilled bytes must not be treated as recovered.\\
273&Prepare an explicit ordering of all four byte positions for this round.\\
275&Visit four offsets starting from the user-selected first byte.\\
276&Wrap the byte position modulo four so every position appears once.\\
278&Recover one complete key byte at a time in the chosen order.\\
279&Mark the byte or group as unconstrained before its first observation; an empty set would instead mean inconsistency.\\
282&Request distinct faults sequentially, numbering observations from one within this target byte/group.\\
283&Query the simulated target for one new faulty ciphertext. No key or teacher trace is returned to this solver.\\
285&Undo the same recovered later rounds in the faulty ciphertext so both outputs refer to the same target round.\\
287&Choose the default known-byte filter unless another compatible filter, such as the Giraud-inspired one, was explicitly supplied.\\
289&Apply the selected-round byte equation to this clean/faulty pair and known mask.\\
290&The if-condition at line 287 was false; follow this alternative path rather than the preceding branch.\\
291&Run the alternate filter and receive its inferred byte position alongside candidate sets and difference.\\
293&Reject an alternate-filter result whose changed byte is inconsistent with the requested target byte.\\
294&Stop this path with the stated exception. The message identifies the violated input, model or verification condition; no recovered key is accepted by this line.\\
296&Distinguish the first observation from subsequent intersections; no previous candidate set exists yet.\\
297&Before a first byte observation all 256 byte values are still possible.\\
298&The first pair initializes the byte set; subsequent pairs intersect it. No new guess can be introduced by intersection.\\
299&The if-condition at line 296 was false; follow this alternative path rather than the preceding branch.\\
300&Record the number of hypotheses remaining before applying the new observation.\\
301&Intersect the prior hypotheses with those from the new pair; never introduce a new candidate.\\
304&An empty intersection means the observations and model cannot be reconciled; do not claim a key.\\
305&Stop this path with the stated exception. The message identifies the violated input, model or verification condition; no recovered key is accepted by this line.\\
307&Count one faulty ciphertext actually consumed by the solver, not a printed stage or a schedule replay.\\
308&Store the surviving count for this particular round and byte.\\
309&Prepare a display-only dictionary of already recovered later subkeys.\\
311&Visit only subkeys already established by the solver.\\
312&Format each recovered subkey as eight hexadecimal digits for event metadata.\\
314&Package the measurement pair, restored views, known mask, input guesses, surviving key guesses and recovered peeling keys for presentation.\\
315&Preserve this observation in order for later inspection and reproducibility.\\
317&Only invoke the optional reporting callback if the caller requested it.\\
319&Pass copied event metadata to reporting. Candidate-set intersection has already been decided in the solver.\\
321&Uniqueness is the acceptance rule: one value remains compatible with all observations for this byte/group.\\
323&Extract the sole surviving byte and store it at the correct big-endian position.\\
324&Stop this observation loop immediately after uniqueness; do not query unused masks.\\
325&The for-loop beginning at line 282 exhausted its observation budget without break; execute the following ambiguity result.\\
328&Return an explicit incomplete result when the observation loop exhausts its budget without a break. No master key is selected.\\
330&Concatenate the four uniquely recovered bytes into one 32-bit round key and make it available for future peeling.\\
333&Prepare the four recovered tail keys in forward key-schedule order.\\
335&Read rk28, rk29, rk30 and rk31 in the order expected by schedule inversion.\\
336&Add the actual recovered subkey to the inverse-schedule seed.\\
338&Arrange recovered rk28..rk31 in schedule order and reverse expansion to derive the original key.\\
340&Return the complete key together with actual recovered tail keys, history, counts and fault-query cost.\\
\bottomrule\end{longtable}\endgroup
\section{\texttt{main}}
Parse the attack-specific command and enter the same SM4 terminal lesson used by run.py.
\begin{center}\begin{tikzpicture}[node distance=4mm]
\node[box,text width=100mm] (input) {Terminal command and options};
\node[box,below=of input,text width=100mm] (action) {Validate options; call sm4\_demo};
\node[box,below=of action,text width=100mm] (output) {Success or ambiguity exit status};
\draw[flow](input)--(action);\draw[flow](action)--(output);
\end{tikzpicture}\end{center}
\remember{\textbf{Physical interpretation:} This is a PC entry point. Selecting --row chooses a simulated byte, not a physical chip coordinate. The compatibility options --column and --seed are parsed but not used in this SM4 path.}
\subsection*{Exact source}
\VerbatimInput[firstnumber=343]{generated/annotated_piret_sm4_main.txt}
\subsection*{Read each executable line}
\begingroup\small\begin{longtable}{r >{\raggedright\arraybackslash}p{138mm}}\toprule Line&Meaning and effect\\\midrule\endfirsthead\toprule Line&Meaning and effect\\\midrule\endhead
343&Declare this callable and its parameters; the body runs only when called. Parse the attack-specific command and enter the same SM4 terminal lesson used by run.py.\\
346&Create the command-line option parser using the module description for help text.\\
348&Register --algorithm. Select the allowed cipher branch for this entry point.\\
349&Register --step. Enable pauses at teaching boundaries.\\
350&Register --color. Select automatic, forced or disabled ANSI colors.\\
351&Register --faults. Set the maximum distinct mask observations per SM4 target byte.\\
352&Register --row. Choose the first simulated SM4 byte position.\\
354&Register --column. Compatibility setting parsed but not used by this SM4 entry point.\\
355&Register --seed. Used for AES mask sampling in run.py; not used for SM4 mask order.\\
357&Decode the supplied argument list (or process command line) and enforce the declared argument types and choices.\\
359&Reject an observation budget outside 1..255 before requesting any faulty outputs.\\
360&Report invalid configuration and terminate argument processing rather than starting an invalid experiment.\\
362&Load the shared terminal controller at call time. This avoids a top-level circular import between attack and presentation modules.\\
364&Run the shared SM4 lesson and test whether full recovery verified.\\
365&Return successful process status after verified complete recovery.\\
367&Return the explicit ambiguity/failure status when full recovery was not verified.\\
\bottomrule\end{longtable}\endgroup
\chapter{Line-by-line: \texttt{run.py}}
\section{\texttt{block}}
Convert a hexadecimal input string to exactly sixteen bytes.
\begin{center}\begin{tikzpicture}[node distance=4mm]
\node[box,text width=100mm] (input) {32 hexadecimal characters};
\node[box,below=of input,text width=100mm] (action) {Decode hex and check length};
\node[box,below=of action,text width=100mm] (output) {16 raw bytes};
\draw[flow](input)--(action);\draw[flow](action)--(output);
\end{tikzpicture}\end{center}
\remember{\textbf{Physical interpretation:} A future host transport would send these bytes to firmware. This function only prepares host memory.}
\subsection*{Exact source}
\VerbatimInput[firstnumber=12]{generated/annotated_run_block.txt}
\subsection*{Read each executable line}
\begingroup\small\begin{longtable}{r >{\raggedright\arraybackslash}p{138mm}}\toprule Line&Meaning and effect\\\midrule\endfirsthead\toprule Line&Meaning and effect\\\midrule\endhead
12&Declare this callable and its parameters; the body runs only when called. Convert a hexadecimal input string to exactly sixteen bytes.\\
13&Decode hexadecimal text into raw bytes; malformed hex raises ValueError instead of silently changing the input.\\
15&Accept only one complete 128-bit block or key for this demonstration.\\
16&Stop this path with the stated exception. The message identifies the violated input, model or verification condition; no recovered key is accepted by this line.\\
18&Return the validated raw bytes used by the cipher and target simulator.\\
\bottomrule\end{longtable}\endgroup
\section{\texttt{LessonTerminal}}
Format the same data as clean/faulty/XOR rows with red differences.
\begin{center}\begin{tikzpicture}[node distance=4mm]
\node[box,text width=100mm] (input) {Two arrays of bytes};
\node[box,below=of input,text width=100mm] (action) {Compare positions; apply ANSI colors};
\node[box,below=of action,text width=100mm] (output) {Readable terminal and log rows};
\draw[flow](input)--(action);\draw[flow](action)--(output);
\end{tikzpicture}\end{center}
\remember{\textbf{Physical interpretation:} Only the PC display changes. Red text is a visualization of differing data, not a measured electrical pulse.}
\subsection*{Exact source}
\VerbatimInput[firstnumber=21]{generated/annotated_run_LessonTerminal.txt}
\subsection*{Read each executable line}
\begingroup\small\begin{longtable}{r >{\raggedright\arraybackslash}p{138mm}}\toprule Line&Meaning and effect\\\midrule\endfirsthead\toprule Line&Meaning and effect\\\midrule\endhead
21&Declare this callable and its parameters; the body runs only when called. Format the same data as clean/faulty/XOR rows with red differences.\\
22&Define the byte-formatting method used to color differences without changing any byte values.\\
23&Accumulate or append the explicitly formatted hexadecimal display fragments; this does not modify cipher data.\\
25&Compare each faulty output byte with its reference position before choosing its color.\\
26&Use red formatting only when this faulty byte differs from the clean reference byte.\\
27&Prepare this changed byte with ANSI red and an immediate reset.\\
28&The if-condition at line 26 was false; follow this alternative path rather than the preceding branch.\\
29&Prepare an unchanged byte as ordinary two-digit hexadecimal text.\\
31&Accumulate or append the explicitly formatted hexadecimal display fragments; this does not modify cipher data.\\
33&Join the prepared hex fragments with spaces for terminal display.\\
35&Define the clean/faulty/XOR printing method used throughout both lessons.\\
36&Accumulate or append the explicitly formatted hexadecimal display fragments; this does not modify cipher data.\\
38&Format each clean reference byte without fault coloring.\\
39&Accumulate or append the explicitly formatted hexadecimal display fragments; this does not modify cipher data.\\
41&Print the text/label "Correct:". This is presentation output, also mirrored to the active log; it performs no target query and changes no candidates.\\
42&Print the text/label "Faulty:". This is presentation output, also mirrored to the active log; it performs no target query and changes no candidates.\\
44&Prepare an explicit byte buffer for the following comparison or recovered-key serialization.\\
46&Visit the aligned correct and faulty bytes in the same serialized order.\\
47&Append their byte XOR, which is zero exactly when the two positions agree.\\
49&Serialize the explicitly constructed difference buffer.\\
51&Print the text/label "XOR:". This is presentation output, also mirrored to the active log; it performs no target query and changes no candidates.\\
\bottomrule\end{longtable}\endgroup
\section{\texttt{aes\_demo}}
Run the AES round-9 lesson when AES is selected; this function is not executed by --algorithm sm4.
\begin{center}\begin{tikzpicture}[node distance=4mm]
\node[box,text width=100mm] (input) {AES selection only};
\node[box,below=of input,text width=100mm] (action) {AES encryption and quartet recovery};
\node[box,below=of action,text width=100mm] (output) {AES verification result};
\draw[flow](input)--(action);\draw[flow](action)--(output);
\end{tikzpicture}\end{center}
\remember{\textbf{Physical interpretation:} Included for complete run.py coverage. It does not contribute to SM4 recovery; all operations here remain host-side software simulation.}
\subsection*{Exact source}
\VerbatimInput[firstnumber=54]{generated/annotated_run_aes_demo.txt}
\subsection*{Read each executable line}
\begingroup\small\begin{longtable}{r >{\raggedright\arraybackslash}p{138mm}}\toprule Line&Meaning and effect\\\midrule\endfirsthead\toprule Line&Meaning and effect\\\midrule\endhead
54&Declare this callable and its parameters; the body runs only when called. Run the AES round-9 lesson when AES is selected; this function is not executed by --algorithm sm4.\\
55&Load AES simulation inputs. This line belongs to the AES branch and does not supply SM4 key material.\\
57&Print the named teaching-section boundary. This affects terminal/log organization only and does not advance the target or the solver.\\
59&Print the text/label "Plaintext:  Key:". This is presentation output, also mirrored to the active log; it performs no target query and changes no candidates.\\
60&Print the text/label "AES layout: index = 4*column + row. This execution has no fault.". This is presentation output, also mirrored to the active log; it performs no target query and changes no candidates.\\
62&Use the AES trace adapter with fault\_round=9; mask zero creates the clean reference.\\
64&Check that the trace adapter matches ordinary AES encryption before teaching its fault behavior.\\
65&Stop this path with the stated exception. The message identifies the violated input, model or verification condition; no recovered key is accepted by this line.\\
67&Compute AES round keys for the teacher-only normal-encryption display.\\
69&Show each AES expanded round key in sequence.\\
70&Print the text/label "K  =". Substitute these runtime fields: number, aes.state\_to\_bytes(round\_key).hex().upper(). This is presentation output, also mirrored to the active log; it performs no target query and changes no candidates.\\
72&Convert the plaintext into the AES column-major matrix for initial key addition.\\
74&Apply the initial AES K0 XOR in the teacher view.\\
76&Print the text/label "Initial AddRoundKey:". This is presentation output, also mirrored to the active log; it performs no target query and changes no candidates.\\
78&Walk the clean AES transformation snapshots in execution order.\\
79&Exclude the adapter marker from the ordinary AES transformation display.\\
80&Print the text/label "R   :". Substitute these runtime fields: number, operation, state.hex().upper(). This is presentation output, also mirrored to the active log; it performs no target query and changes no candidates.\\
81&Wait for Enter only if --step is enabled. Pausing does not change a state word, create a fault or alter a candidate set.\\
83&Print the text/label "Correct ciphertext:". This is presentation output, also mirrored to the active log; it performs no target query and changes no candidates.\\
85&Print the named teaching-section boundary. This affects terminal/log organization only and does not advance the target or the solver.\\
87&Print the text/label "Recover one K10 quartet per column; repeat across all four columns.". This is presentation output, also mirrored to the active log; it performs no target query and changes no candidates.\\
88&Print the text/label "d = InvSbox[C[p] XOR k] XOR InvSbox[F[p] XOR k].". This is presentation output, also mirrored to the active log; it performs no target query and changes no candidates.\\
89&Print the text/label "Four d values must match one of 4*255 = 1020 MixColumns difference patterns.". This is presentation output, also mirrored to the active log; it performs no target query and changes no candidates.\\
91&Use the two preferred initial AES masks when the observation budget allows both.\\
92&Try AES masks 1 and 2 first, then reproducibly sample other distinct nonzero masks up to the limit.\\
93&The if-condition at line 91 was false; follow this alternative path rather than the preceding branch.\\
94&For a budget of one, run only the first AES mask rather than requesting extra observations.\\
96&Prepare sixteen AES final-key bytes and a total faulty-output counter.\\
98&Recover each of the four AES pre-ShiftRows key quartets in turn.\\
99&Mark the byte or group as unconstrained before its first observation; an empty set would instead mean inconsistency.\\
101&Print the text/label "COLUMN  : output/key positions". Substitute these runtime fields: column, piret\_aes.output\_positions(column). This is presentation output, also mirrored to the active log; it performs no target query and changes no candidates.\\
103&Request distinct faults sequentially, numbering observations from one within this target byte/group.\\
104&Run an AES round-9 byte-fault encryption for the selected column and mask.\\
105&Count this faulty AES encryption once, independently of display stages.\\
106&Print the text/label "Observation  : simulator XOR  , row 0 / column  .". Substitute these runtime fields: observation, mask, column. This is presentation output, also mirrored to the active log; it performs no target query and changes no candidates.\\
107&Print the text/label "The solver enumerates unknown fault rows/values and does not receive this mask.". This is presentation output, also mirrored to the active log; it performs no target query and changes no candidates.\\
108&Index clean AES snapshots by round and operation for paired fault-propagation comparisons.\\
110&Build the clean AES snapshot lookup explicitly from each recorded round and operation.\\
111&Index the clean state by its round/operation pair for later fault-propagation comparisons.\\
113&Visit each faulty AES transformation snapshot for the current observation.\\
114&Display only the AES fault round and later propagation in this part of the lesson.\\
115&Print the text/label "R". Substitute these runtime fields: number, operation. This is presentation output, also mirrored to the active log; it performs no target query and changes no candidates.\\
116&Compare the two supplied byte arrays visually, compute their XOR and color changed positions red. The data is not modified.\\
117&Wait for Enter only if --step is enabled. Pausing does not change a state word, create a fault or alter a candidate set.\\
119&Compare the two supplied byte arrays visually, compute their XOR and color changed positions red. The data is not modified.\\
120&Print exhaustive AES byte guesses and shared-error quartet patterns; this branch is not used in SM4-only execution.\\
122&Derive compatible AES K10 quartets from correct/faulty ciphertexts, not from the printed key.\\
124&Distinguish the first observation from subsequent intersections; no previous candidate set exists yet.\\
125&Mark the initial AES quartet history as unconstrained rather than claiming a previous finite set.\\
126&The if-condition at line 124 was false; follow this alternative path rather than the preceding branch.\\
127&Record the number of hypotheses remaining before applying the new observation.\\
130&Distinguish the first observation from subsequent intersections; no previous candidate set exists yet.\\
131&Initialize this AES quartet set directly from the first observed pair.\\
132&The if-condition at line 130 was false; follow this alternative path rather than the preceding branch.\\
133&Intersect the prior hypotheses with those from the new pair; never introduce a new candidate.\\
135&Print the text/label "Candidates in this pair:  ; intersection:   ->". Substitute these runtime fields: len(fresh), before, len(survivors). This is presentation output, also mirrored to the active log; it performs no target query and changes no candidates.\\
136&Accumulate or append the explicitly formatted hexadecimal display fragments; this does not modify cipher data.\\
138&Format up to four candidates for the short summary; the earlier detailed tables remain exhaustive.\\
139&Accumulate or append the explicitly formatted hexadecimal display fragments; this does not modify cipher data.\\
141&Print the text/label "First survivors:   |". This is presentation output, also mirrored to the active log; it performs no target query and changes no candidates.\\
142&Wait for Enter only if --step is enabled. Pausing does not change a state word, create a fault or alter a candidate set.\\
144&An empty intersection means the observations and model cannot be reconciled; do not claim a key.\\
145&Stop this path with the stated exception. The message identifies the violated input, model or verification condition; no recovered key is accepted by this line.\\
147&Uniqueness is the acceptance rule: one value remains compatible with all observations for this byte/group.\\
148&Place the unique AES tuple bytes into final ciphertext/key order using inverse ShiftRows mapping.\\
149&Store one uniquely recovered AES final-key byte at its true serialized position.\\
151&Stop this observation loop immediately after uniqueness; do not query unused masks.\\
152&The for-loop beginning at line 103 exhausted its observation budget without break; execute the following ambiguity result.\\
153&Print the text/label "AES recovery remains ambiguous. Increase --observations. 31". This is presentation output, also mirrored to the active log; it performs no target query and changes no candidates.\\
154&Tell the caller that this requested lesson did not establish verified complete recovery.\\
156&Print the named teaching-section boundary. This affects terminal/log organization only and does not advance the target or the solver.\\
158&Print the text/label "Recovered K10:". This is presentation output, also mirrored to the active log; it performs no target query and changes no candidates.\\
160&Derive AES K0 from the recovered K10 through the detailed display helper.\\
162&Run and print recovered-key encryption/decryption consistency checks on the ordinary cipher.\\
164&Check the AES recovered key only after recovery: ground truth and re-encryption must both agree.\\
165&Stop this path with the stated exception. The message identifies the violated input, model or verification condition; no recovered key is accepted by this line.\\
167&Print the text/label "32 AES MASTER KEY:   / PASS (  faulty outputs)". Substitute these runtime fields: recovered.hex().upper(), total. This is presentation output, also mirrored to the active log; it performs no target query and changes no candidates.\\
169&Return verified success to the command-line controller.\\
\bottomrule\end{longtable}\endgroup
\section{\texttt{inverse\_rows}}
Delegate inverse-matrix construction to the SM4 attack module.
\begin{center}\begin{tikzpicture}[node distance=4mm]
\node[box,text width=100mm] (input) {Public linear map};
\node[box,below=of input,text width=100mm] (action) {Call piret\_sm4.inverse\_rows};
\node[box,below=of action,text width=100mm] (output) {Cached inverse-matrix rows};
\draw[flow](input)--(action);\draw[flow](action)--(output);
\end{tikzpicture}\end{center}
\remember{\textbf{Physical interpretation:} This wrapper exists on the host for compatible call sites; it does not issue a hardware operation.}
\subsection*{Exact source}
\VerbatimInput[firstnumber=172]{generated/annotated_run_inverse_rows.txt}
\subsection*{Read each executable line}
\begingroup\small\begin{longtable}{r >{\raggedright\arraybackslash}p{138mm}}\toprule Line&Meaning and effect\\\midrule\endfirsthead\toprule Line&Meaning and effect\\\midrule\endhead
172&Declare this callable and its parameters; the body runs only when called. Delegate inverse-matrix construction to the SM4 attack module.\\
175&Delegate to the single cached implementation of inverse-matrix construction.\\
\bottomrule\end{longtable}\endgroup
\section{\texttt{inverse\_linear}}
Delegate inverse diffusion to the SM4 attack module.
\begin{center}\begin{tikzpicture}[node distance=4mm]
\node[box,text width=100mm] (input) {Observed difference};
\node[box,below=of input,text width=100mm] (action) {Call piret\_sm4.inverse\_linear};
\node[box,below=of action,text width=100mm] (output) {Restored S-box difference};
\draw[flow](input)--(action);\draw[flow](action)--(output);
\end{tikzpicture}\end{center}
\remember{\textbf{Physical interpretation:} Offline PC algebra; no new target encryption occurs.}
\subsection*{Exact source}
\VerbatimInput[firstnumber=178]{generated/annotated_run_inverse_linear.txt}
\subsection*{Read each executable line}
\begingroup\small\begin{longtable}{r >{\raggedright\arraybackslash}p{138mm}}\toprule Line&Meaning and effect\\\midrule\endfirsthead\toprule Line&Meaning and effect\\\midrule\endhead
178&Declare this callable and its parameters; the body runs only when called. Delegate inverse diffusion to the SM4 attack module.\\
181&Delegate inverse diffusion; the wrapper introduces no new equation or target operation.\\
\bottomrule\end{longtable}\endgroup
\section{\texttt{sm4\_encrypt}}
Delegate a clean or faulty encryption request to the software target adapter.
\begin{center}\begin{tikzpicture}[node distance=4mm]
\node[box,text width=100mm] (input) {Plaintext, key and optional fault};
\node[box,below=of input,text width=100mm] (action) {Call SM4 target adapter};
\node[box,below=of action,text width=100mm] (output) {Ciphertext and trace};
\draw[flow](input)--(action);\draw[flow](action)--(output);
\end{tikzpicture}\end{center}
\remember{\textbf{Physical interpretation:} This wrapper currently calls Python. A serial transport is not implemented by this function.}
\subsection*{Exact source}
\VerbatimInput[firstnumber=184]{generated/annotated_run_sm4_encrypt.txt}
\subsection*{Read each executable line}
\begingroup\small\begin{longtable}{r >{\raggedright\arraybackslash}p{138mm}}\toprule Line&Meaning and effect\\\midrule\endfirsthead\toprule Line&Meaning and effect\\\midrule\endhead
184&Declare this callable and its parameters; the body runs only when called. Delegate a clean or faulty encryption request to the software target adapter.\\
187&Forward the target inputs and optional fault descriptor to the SM4 simulator, returning both output and trace.\\
\bottomrule\end{longtable}\endgroup
\section{\texttt{sm4\_byte\_candidates}}
Delegate the known-mask byte equation to the attack module.
\begin{center}\begin{tikzpicture}[node distance=4mm]
\node[box,text width=100mm] (input) {Restored output pair};
\node[box,below=of input,text width=100mm] (action) {Call SM4 candidate solver};
\node[box,below=of action,text width=100mm] (output) {Input/key guesses and difference};
\draw[flow](input)--(action);\draw[flow](action)--(output);
\end{tikzpicture}\end{center}
\remember{\textbf{Physical interpretation:} This is PC analysis of outputs, not access to hidden chip state.}
\subsection*{Exact source}
\VerbatimInput[firstnumber=190]{generated/annotated_run_sm4_byte_candidates.txt}
\subsection*{Read each executable line}
\begingroup\small\begin{longtable}{r >{\raggedright\arraybackslash}p{138mm}}\toprule Line&Meaning and effect\\\midrule\endfirsthead\toprule Line&Meaning and effect\\\midrule\endhead
190&Declare this callable and its parameters; the body runs only when called. Delegate the known-mask byte equation to the attack module.\\
193&Delegate ciphertext-based candidate filtering rather than duplicate it in presentation code.\\
\bottomrule\end{longtable}\endgroup
\section{\texttt{reverse\_sm4}}
Delegate schedule inversion without duplicating its equations.
\begin{center}\begin{tikzpicture}[node distance=4mm]
\node[box,text width=100mm] (input) {Four recovered tail keys};
\node[box,below=of input,text width=100mm] (action) {Call inverse key expansion};
\node[box,below=of action,text width=100mm] (output) {Original key bytes};
\draw[flow](input)--(action);\draw[flow](action)--(output);
\end{tikzpicture}\end{center}
\remember{\textbf{Physical interpretation:} The host derives key material mathematically after the fault observations.}
\subsection*{Exact source}
\VerbatimInput[firstnumber=196]{generated/annotated_run_reverse_sm4.txt}
\subsection*{Read each executable line}
\begingroup\small\begin{longtable}{r >{\raggedright\arraybackslash}p{138mm}}\toprule Line&Meaning and effect\\\midrule\endfirsthead\toprule Line&Meaning and effect\\\midrule\endhead
196&Declare this callable and its parameters; the body runs only when called. Delegate schedule inversion without duplicating its equations.\\
199&Delegate schedule reversal, forwarding the optional explanatory callback.\\
\bottomrule\end{longtable}\endgroup
\section{\texttt{show\_sm4\_peeling}}
Explain every reverse round numerically for a correct or faulty output.
\begin{center}\begin{tikzpicture}[node distance=4mm]
\node[box,text width=100mm] (input) {Ciphertext and recovered later keys};
\node[box,below=of input,text width=100mm] (action) {Print each reconstructed older word};
\node[box,below=of action,text width=100mm] (output) {Detailed peeling explanation};
\draw[flow](input)--(action);\draw[flow](action)--(output);
\end{tikzpicture}\end{center}
\remember{\textbf{Physical interpretation:} This reconstructs a software view on the PC. It is not a live trace of past target registers.}
\subsection*{Exact source}
\VerbatimInput[firstnumber=202]{generated/annotated_run_show_sm4_peeling.txt}
\subsection*{Read each executable line}
\begingroup\small\begin{longtable}{r >{\raggedright\arraybackslash}p{138mm}}\toprule Line&Meaning and effect\\\midrule\endfirsthead\toprule Line&Meaning and effect\\\midrule\endhead
202&Declare this callable and its parameters; the body runs only when called. Explain every reverse round numerically for a correct or faulty output.\\
205&Start the explanatory peeling pass from the measured final ciphertext.\\
208&Undo only rounds later than the requested target round, newest first.\\
209&Restore chronological word order for the current reverse-round explanation.\\
211&Use the already recovered subkey corresponding to the later round being undone.\\
213&Form that round's S-box input from the three newest surviving words and its recovered key.\\
215&Substitute the four reconstructed input bytes using the public SM4 S-box.\\
217&Apply encryption diffusion L to the substituted value for the reverse-round equation.\\
220&Ask the actual attack helper for the restored view, so the explanatory pass cannot silently use different reversal logic.\\
222&Read the newly reconstructed oldest word from the last word of the ciphertext-style restored view.\\
224&Print the text/label "Undo round   with recovered rk =". Substitute these runtime fields: number, number, round\_key. This is presentation output, also mirrored to the active log; it performs no target query and changes no candidates.\\
225&Print the text/label "A =   XOR   XOR   XOR   =". Substitute these runtime fields: words[0], words[1], words[2], round\_key, value. This is presentation output, also mirrored to the active log; it performs no target query and changes no candidates.\\
226&Print the text/label "tau(A) =  ; L(tau(A)) =". Substitute these runtime fields: substituted, diffused. This is presentation output, also mirrored to the active log; it performs no target query and changes no candidates.\\
227&Print the text/label "X  = X  XOR L(tau(A)) =   XOR   =". Substitute these runtime fields: number, number + 4, words[3], diffused, previous. This is presentation output, also mirrored to the active log; it performs no target query and changes no candidates.\\
228&Print the text/label "Restored output view:". Substitute these runtime fields: view.hex().upper(). This is presentation output, also mirrored to the active log; it performs no target query and changes no candidates.\\
\bottomrule\end{longtable}\endgroup
\section{\texttt{show\_sm4\_key\_schedule}}
Produce and print detailed SM4 inverse-schedule teaching events.
\begin{center}\begin{tikzpicture}[node distance=4mm]
\node[box,text width=100mm] (input) {DFA-recovered tail keys};
\node[box,below=of input,text width=100mm] (action) {Explain all reverse words and FK removal};
\node[box,below=of action,text width=100mm] (output) {Master key and printed derivation};
\draw[flow](input)--(action);\draw[flow](action)--(output);
\end{tikzpicture}\end{center}
\remember{\textbf{Physical interpretation:} This display may repeat an already completed inverse schedule. That repetition does not collect additional faults.}
\subsection*{Exact source}
\VerbatimInput[firstnumber=231]{generated/annotated_run_show_sm4_key_schedule.txt}
\subsection*{Read each executable line}
\begingroup\small\begin{longtable}{r >{\raggedright\arraybackslash}p{138mm}}\toprule Line&Meaning and effect\\\midrule\endfirsthead\toprule Line&Meaning and effect\\\midrule\endhead
231&Declare this callable and its parameters; the body runs only when called. Produce and print detailed SM4 inverse-schedule teaching events.\\
234&Build all numerical inverse-schedule events and the resulting SM4 master key from recovered tail keys.\\
236&Print the numerical teaching events already generated by the schedule helper. This is explanatory replay, not additional fault collection.\\
238&Return the master-key bytes produced by the explained inverse schedule.\\
\bottomrule\end{longtable}\endgroup
\section{\texttt{show\_sm4\_observation}}
Explain one observation, its inverse-L calculation, its candidates and its uniqueness status.
\begin{center}\begin{tikzpicture}[node distance=4mm]
\node[box,text width=100mm] (input) {Solver event and previous sets};
\node[box,below=of input,text width=100mm] (action) {Render equations and intersections};
\node[box,below=of action,text width=100mm] (output) {Recovered-byte/round-key summaries};
\draw[flow](input)--(action);\draw[flow](action)--(output);
\end{tikzpicture}\end{center}
\remember{\textbf{Physical interpretation:} The event contains measurements and reconstructed values. Teacher traces and hidden states must not be confused with the solver inputs.}
\subsection*{Exact source}
\VerbatimInput[firstnumber=241]{generated/annotated_run_show_sm4_observation.txt}
\subsection*{Read each executable line}
\begingroup\small\begin{longtable}{r >{\raggedright\arraybackslash}p{138mm}}\toprule Line&Meaning and effect\\\midrule\endfirsthead\toprule Line&Meaning and effect\\\midrule\endhead
241&Declare this callable and its parameters; the body runs only when called. Explain one observation, its inverse-L calculation, its candidates and its uniqueness status.\\
244&Read the selected target round and byte from the completed solver observation.\\
247&Use a compound dictionary key so candidate history for one byte cannot contaminate another round or position.\\
249&Print every detailed stage for this observation: restored views, inverse-L contributions, candidate tests and intersection. The previous display set is passed only for explanation.\\
251&Remember the current candidate set for the next detailed intersection display; this does not change the solver.\\
253&Print the round-level introduction only for the first observation of the first displayed byte.\\
254&Print the named teaching-section boundary. This affects terminal/log organization only and does not advance the target or the solver.\\
256&The actual last round requires no peeling; earlier target rounds must remove later rounds first.\\
257&Print the text/label "Round 31 is already last; no round needs to be undone.". This is presentation output, also mirrored to the active log; it performs no target query and changes no candidates.\\
258&The if-condition at line 256 was false; follow this alternative path rather than the preceding branch.\\
259&Prepare the decoded recovered-key dictionary for the explanatory peeling pass.\\
261&Decode each already recovered later key carried by the event.\\
262&Convert string indices and hexadecimal key labels back to integers for the display calculation.\\
264&Print the text/label "Restore the target round using ONLY previously recovered later subkeys.". This is presentation output, also mirrored to the active log; it performs no target query and changes no candidates.\\
265&Print the text/label "Correct execution:". This is presentation output, also mirrored to the active log; it performs no target query and changes no candidates.\\
266&Explain correct or faulty round peeling separately; these reconstructed views remain offline host values.\\
267&Print the text/label "Faulty execution (a NEW fault at the target round):". This is presentation output, also mirrored to the active log; it performs no target query and changes no candidates.\\
268&Explain correct or faulty round peeling separately; these reconstructed views remain offline host values.\\
270&Print the text/label "Round  , byte  , observation  , mask  round byte observation mask". Substitute these runtime fields: event['round'], event['byte'], event['observation'], event['mask']. This is presentation output, also mirrored to the active log; it performs no target query and changes no candidates.\\
272&Compare the two supplied byte arrays visually, compute their XOR and color changed positions red. The data is not modified.\\
274&Print the text/label "Undo later rounds using ONLY previously recovered subkeys:  peeled\_with". This is presentation output, also mirrored to the active log; it performs no target query and changes no candidates.\\
276&Compare the two supplied byte arrays visually, compute their XOR and color changed positions red. The data is not modified.\\
279&Decode the clean restored output view into four big-endian words for the explanatory equations.\\
281&Decode the corresponding faulty restored output with exactly the same layout.\\
283&Map byte 0,1,2,3 to bit shifts 24,16,8,0 in a big-endian 32-bit word.\\
285&Recompute q from public restored words for the displayed key-byte conversion.\\
287&Select the target byte of the already computed S-box output difference.\\
289&Print the text/label "DeltaX =   XOR   =". Substitute these runtime fields: clean\_words[0], faulty\_words[0], clean\_words[0] \textasciicircum{} faulty\_words[0]. This is presentation output, also mirrored to the active log; it performs no target query and changes no candidates.\\
290&Print the text/label "DeltaB = L-inverse(DeltaX) =  ; selected byte difference =  delta\_b". Substitute these runtime fields: event['delta\_b'], delta. This is presentation output, also mirrored to the active log; it performs no target query and changes no candidates.\\
291&Print the text/label "Known byte q = byte   of (X  XOR X  XOR X ) =". Substitute these runtime fields: byte, number + 1, number + 2, number + 3, known. This is presentation output, also mirrored to the active log; it performs no target query and changes no candidates.\\
294&Unknown-bit Giraud must test every one-bit mask. The default controlled-byte demonstration instead knows its mask.\\
295&List the eight possible single-bit errors without revealing which simulator bit was actually selected.\\
297&Print the text/label "UNKNOWN-BIT FILTER: the simulator mask above is NOT given to the candidate solver.". This is presentation output, also mirrored to the active log; it performs no target query and changes no candidates.\\
298&The if-condition at line 294 was false; follow this alternative path rather than the preceding branch.\\
299&Use the exact known error for the controlled-byte or explicitly known-bit model.\\
301&Print the text/label "KNOWN-MASK FILTER: the exact error is given to the candidate solver.". This is presentation output, also mirrored to the active log; it performs no target query and changes no candidates.\\
303&Accumulate or append the explicitly formatted hexadecimal display fragments; this does not modify cipher data.\\
305&Visit each value of the displayed allowed-error, survivor or recovered-byte sequence explicitly.\\
306&Accumulate or append the explicitly formatted hexadecimal display fragments; this does not modify cipher data.\\
308&Print the text/label "Allowed input errors:". This is presentation output, also mirrored to the active log; it performs no target query and changes no candidates.\\
309&Print the text/label "Test a=00..FF: at least one allowed e must satisfy S(a) XOR S(a XOR e) =  .". Substitute these runtime fields: delta. This is presentation output, also mirrored to the active log; it performs no target query and changes no candidates.\\
311&Explain every compatible S-box input that survived the current pair equation.\\
312&Start without a compatible error witness; the explicit search must find one.\\
314&Test each error permitted by the selected known-mask or unknown-bit model.\\
315&Construct the corresponding faulty S-box input for this candidate and error.\\
316&Calculate the predicted S-box difference for this error witness.\\
318&Accept the error witness only if it matches the observed byte difference.\\
319&Remember this compatible error before stopping the witness search.\\
320&Stop the permitted-error search once one compatible witness is found.\\
322&Report the impossible-witness condition rather than inventing an error value.\\
323&Stop this path with the stated exception. The message identifies the violated input, model or verification condition; no recovered key is accepted by this line.\\
325&Reconstruct the corresponding faulty S-box input for this compatible hypothesis.\\
327&Print the text/label "a= , compatible e= : S[ ] XOR S[ ] =   XOR   =  ; rk byte = a XOR q =   XOR   =". Substitute these runtime fields: candidate, error, candidate, changed, sm4.SBOX[candidate], sm4.SBOX[changed], delta, candidate, known, candidate \textasciicircum{} known. This is presentation output, also mirrored to the active log; it performs no target query and changes no candidates.\\
329&Print the text/label "L-inverse difference:  ; candidates   ->  delta\_b before remaining". Substitute these runtime fields: event['delta\_b'], event['before'], event['remaining']. This is presentation output, also mirrored to the active log; it performs no target query and changes no candidates.\\
331&Accumulate or append the explicitly formatted hexadecimal display fragments; this does not modify cipher data.\\
333&Visit each value of the displayed allowed-error, survivor or recovered-byte sequence explicitly.\\
334&Accumulate or append the explicitly formatted hexadecimal display fragments; this does not modify cipher data.\\
336&Print the text/label "Survivors:". This is presentation output, also mirrored to the active log; it performs no target query and changes no candidates.\\
338&Only print a recovered-byte declaration when the accumulated intersection has exactly one value.\\
339&Read that unique key-byte value from the sorted event list.\\
342&Maintain a per-round collection of uniquely recovered bytes for final round-key assembly.\\
344&Print the text/label "32 RECOVERED rk  byte   =   / UNIQUE". Substitute these runtime fields: number, byte, value. This is presentation output, also mirrored to the active log; it performs no target query and changes no candidates.\\
346&Wait until all four byte positions are known before claiming a complete 32-bit subkey.\\
347&Prepare an explicit byte buffer for the following comparison or recovered-key serialization.\\
349&Visit the four recovered byte positions in big-endian order.\\
350&Add the uniquely recovered byte at this position to the round-key buffer.\\
352&Serialize all four recovered bytes into the assembled round-key view.\\
353&Accumulate or append the explicitly formatted hexadecimal display fragments; this does not modify cipher data.\\
355&Visit each value of the displayed allowed-error, survivor or recovered-byte sequence explicitly.\\
356&Accumulate or append the explicitly formatted hexadecimal display fragments; this does not modify cipher data.\\
358&Print the text/label "Assemble bytes in big-endian order:   ||". This is presentation output, also mirrored to the active log; it performs no target query and changes no candidates.\\
359&Print the text/label "1;32 DFA ASSEMBLED rk  =". Substitute these runtime fields: number, ordered.hex().upper(). This is presentation output, also mirrored to the active log; it performs no target query and changes no candidates.\\
361&Wait for Enter only if --step is enabled. Pausing does not change a state word, create a fault or alter a candidate set.\\
\bottomrule\end{longtable}\endgroup
\section{\texttt{show\_aes\_key\_schedule}}
Print the AES inverse schedule; this path is inactive in the SM4 demonstration.
\begin{center}\begin{tikzpicture}[node distance=4mm]
\node[box,text width=100mm] (input) {Recovered AES K10};
\node[box,below=of input,text width=100mm] (action) {Derive AES schedule events};
\node[box,below=of action,text width=100mm] (output) {AES K0};
\draw[flow](input)--(action);\draw[flow](action)--(output);
\end{tikzpicture}\end{center}
\remember{\textbf{Physical interpretation:} Included only because it is a function in run.py. The SM4 path instead uses show\_sm4\_key\_schedule.}
\subsection*{Exact source}
\VerbatimInput[firstnumber=364]{generated/annotated_run_show_aes_key_schedule.txt}
\subsection*{Read each executable line}
\begingroup\small\begin{longtable}{r >{\raggedright\arraybackslash}p{138mm}}\toprule Line&Meaning and effect\\\midrule\endfirsthead\toprule Line&Meaning and effect\\\midrule\endhead
364&Declare this callable and its parameters; the body runs only when called. Print the AES inverse schedule; this path is inactive in the SM4 demonstration.\\
367&Build AES-specific inverse-schedule teaching events; this function is inactive for SM4-only execution.\\
369&Print the numerical teaching events already generated by the schedule helper. This is explanatory replay, not additional fault collection.\\
371&Return the master-key bytes produced by the explained inverse schedule.\\
\bottomrule\end{longtable}\endgroup
\section{\texttt{sm4\_demo}}
Build the clean reference, define target/report callbacks, recover the key and validate the result.
\begin{center}\begin{tikzpicture}[node distance=4mm]
\node[box,text width=100mm] (input) {Configured SM4 inputs};
\node[box,below=of input,text width=100mm] (action) {Clean run; faults; recovery; checks};
\node[box,below=of action,text width=100mm] (output) {Printed SM4 result and boolean success};
\draw[flow](input)--(action);\draw[flow](action)--(output);
\end{tikzpicture}\end{center}
\remember{\textbf{Physical interpretation:} This is the host experiment controller. In a board experiment the oracle would need a real transport and glitch campaign instead of the present Python target adapter.}
\subsection*{Exact source}
\VerbatimInput[firstnumber=374]{generated/annotated_run_sm4_demo.txt}
\subsection*{Read each executable line}
\begingroup\small\begin{longtable}{r >{\raggedright\arraybackslash}p{138mm}}\toprule Line&Meaning and effect\\\midrule\endfirsthead\toprule Line&Meaning and effect\\\midrule\endhead
374&Declare this callable and its parameters; the body runs only when called. Build the clean reference, define target/report callbacks, recover the key and validate the result.\\
377&Read and validate the configured target key and plaintext for simulation. The solver does not receive the key.\\
379&Run a clean reference encryption with no fault selected and retain teacher-only round records.\\
381&Print the named teaching-section boundary. This affects terminal/log organization only and does not advance the target or the solver.\\
383&Print each of the thirty-two clean encryption-round records in order.\\
384&Print the text/label "Round  : A=  B=  L(B)=  new= round a b linear new". Substitute these runtime fields: item['round'], item['a'], item['b'], item['linear'], item['new']. This is presentation output, also mirrored to the active log; it performs no target query and changes no candidates.\\
386&Wait for Enter only if --step is enabled. Pausing does not change a state word, create a fault or alter a candidate set.\\
388&Print the text/label "Ciphertext:". This is presentation output, also mirrored to the active log; it performs no target query and changes no candidates.\\
390&Print the named teaching-section boundary. This affects terminal/log organization only and does not advance the target or the solver.\\
392&Print the text/label "Controlled faults in EACH final round are required. No reference subkeys are supplied.". This is presentation output, also mirrored to the active log; it performs no target query and changes no candidates.\\
393&Print the text/label "For each byte: undo L, test S(a) XOR S(a XOR mask), then convert a to a subkey guess.". This is presentation output, also mirrored to the active log; it performs no target query and changes no candidates.\\
394&Print the text/label "Repeat independent faults and intersect guesses. Only one survivor identifies a byte.". This is presentation output, also mirrored to the active log; it performs no target query and changes no candidates.\\
396&Start presentation-only histories for candidate sets and uniquely recovered round-key bytes.\\
399&Define a target-query callback closing over the plaintext and simulated secret key. Its public result is only a ciphertext.\\
400&Perform one new target encryption with the selected fault. The detailed trace is for teaching, not solving.\\
402&Print teacher-only fault propagation from the target trace. This exposes intermediate values for learning, but none are returned as solver inputs.\\
404&Give the attacker only the target ciphertext; do not return its trace or target key.\\
406&Define the presentation callback invoked after each solver observation; it does not generate faults or choose candidates.\\
407&Render a solver event without taking control of its candidate computation.\\
409&Run full output-based SM4 recovery with mask budget 1..limit, selected starting byte and a reporting callback.\\
411&Detect an observation budget exhausted before full uniqueness; do not continue to a claimed master key.\\
412&Print the text/label "SM4 recovery remains ambiguous. Increase --observations.". This is presentation output, also mirrored to the active log; it performs no target query and changes no candidates.\\
414&Tell the caller that this requested lesson did not establish verified complete recovery.\\
416&Summarize only the four subkeys actually recovered through fault observations.\\
417&Print the text/label "32 DFA RECOVERED rk  =   / PASS". Substitute these runtime fields: number, value. This is presentation output, also mirrored to the active log; it performs no target query and changes no candidates.\\
419&Repeat inverse expansion through the detailed teaching view and retain its reconstructed master key.\\
421&Cross-check the detailed display's inverse schedule against the recovery engine's result.\\
422&Stop this path with the stated exception. The message identifies the violated input, model or verification condition; no recovered key is accepted by this line.\\
424&Run and print recovered-key encryption/decryption consistency checks on the ordinary cipher.\\
426&Re-encrypt the known plaintext with the recovered key using the ordinary SM4 cipher.\\
428&Validate both simulation ground-truth equality and the original correct ciphertext, only after recovery.\\
429&Stop this path with the stated exception. The message identifies the violated input, model or verification condition; no recovered key is accepted by this line.\\
431&Print the named teaching-section boundary. This affects terminal/log organization only and does not advance the target or the solver.\\
433&Print the text/label "Plaintext:". This is presentation output, also mirrored to the active log; it performs no target query and changes no candidates.\\
434&Print the text/label "Expected ciphertext:". This is presentation output, also mirrored to the active log; it performs no target query and changes no candidates.\\
435&Print the text/label "Recovered-key ciphertext:". This is presentation output, also mirrored to the active log; it performs no target query and changes no candidates.\\
436&Print the text/label "32  / PASS SM4 MASTER KEY:". This is presentation output, also mirrored to the active log; it performs no target query and changes no candidates.\\
437&Print the text/label "faulty outputs. All four subkeys came from DFA, not references. used". Substitute these runtime fields: result['used']. This is presentation output, also mirrored to the active log; it performs no target query and changes no candidates.\\
439&Return verified success to the command-line controller.\\
\bottomrule\end{longtable}\endgroup
\section{\texttt{main}}
Choose AES, SM4 or both and aggregate their exit status.
\begin{center}\begin{tikzpicture}[node distance=4mm]
\node[box,text width=100mm] (input) {Command-line arguments};
\node[box,below=of input,text width=100mm] (action) {Validate configuration; select lessons};
\node[box,below=of action,text width=100mm] (output) {Process success/ambiguity code};
\draw[flow](input)--(action);\draw[flow](action)--(output);
\end{tikzpicture}\end{center}
\remember{\textbf{Physical interpretation:} All selection, printing and log handling runs on the PC. --algorithm sm4 skips the AES lesson, although configuration validation checks both modules.}
\subsection*{Exact source}
\VerbatimInput[firstnumber=442]{generated/annotated_run_main.txt}
\subsection*{Read each executable line}
\begingroup\small\begin{longtable}{r >{\raggedright\arraybackslash}p{138mm}}\toprule Line&Meaning and effect\\\midrule\endfirsthead\toprule Line&Meaning and effect\\\midrule\endhead
442&Declare this callable and its parameters; the body runs only when called. Choose AES, SM4 or both and aggregate their exit status.\\
443&Create the command-line option parser using the module description for help text.\\
445&Register --algorithm. Select the allowed cipher branch for this entry point.\\
446&Register --step. Enable pauses at teaching boundaries.\\
447&Register --color. Select automatic, forced or disabled ANSI colors.\\
448&Register --observations. Set the per-byte or per-column observation budget.\\
449&Register --seed. Used for AES mask sampling in run.py; not used for SM4 mask order.\\
451&Decode the supplied argument list (or process command line) and enforce the declared argument types and choices.\\
453&Reject an observation budget outside 1..255 before requesting any faulty outputs.\\
454&Report invalid configuration and terminate argument processing rather than starting an invalid experiment.\\
456&Begin a protected configuration-validation block; the associated except path turns invalid input into a parser error.\\
457&Validate both cipher configurations even if just one lesson was requested; this is configuration checking, not AES execution.\\
458&Require this configured key or plaintext to decode as a complete block.\\
459&Catch the named validation exception and report its message through the command-line parser.\\
460&Report invalid configuration and terminate argument processing rather than starting an invalid experiment.\\
462&Configure presentation colors and optional pauses; these choices do not affect the cryptographic equations.\\
464&Print the text/label "AES\_SM4\_DFA.pdf / chapters 2-6 / software demonstration". This is presentation output, also mirrored to the active log; it performs no target query and changes no candidates.\\
465&Print the text/label "RED marks bytes changed by the injected fault, including propagated changes.". This is presentation output, also mirrored to the active log; it performs no target query and changes no candidates.\\
468&Collect independent success flags for the requested algorithm demonstrations.\\
470&Execute the AES branch only when requested. It is skipped by --algorithm sm4.\\
471&Run the AES lesson and store its success flag; not executed in the SM4-only command.\\
473&Execute the SM4 demonstration for either an explicit SM4 selection or the combined default.\\
474&Run the SM4 lesson with its observation budget and store whether full recovery verified.\\
476&Check the already collected success flags for all requested algorithm demonstrations.\\
477&Select process success because all requested demonstrations verified.\\
478&The if-condition at line 476 was false; follow this alternative path rather than the preceding branch.\\
479&Select process ambiguity status because a requested demonstration did not fully recover its key.\\
481&Return the word assembled from all thirty-two inverse-row parity calculations.\\
\bottomrule\end{longtable}\endgroup